\chapter{Conclusion}

This thesis examined SVG generation and vectorization methods from the perspective of representation, supervision, and structural quality. While recent progress in generative modeling has significantly advanced raster image synthesis, vector graphics generation introduces additional challenges due to the structured nature of SVG. Unlike raster outputs, SVG graphics must preserve geometry, hierarchy, and editability, which makes evaluation and model design fundamentally different from conventional image generation tasks.

Through a comparative analysis of representative models, this work showed that the internal representation used by a method strongly influences the structure of the resulting SVG. Parametric optimization approaches provide direct geometric control but often produce redundant paths when supervision operates primarily in raster space. Sequential and dataset-trained models can generate coherent structure within constrained domains, while implicit and layered neural representations introduce architectural constraints that encourage compact and interpretable decomposition. These observations indicate that structural quality is not solely an optimization outcome but emerges from representational design.

The analysis further demonstrated that supervision signals shape geometric behavior in complementary ways. Pixel-based reconstruction encourages accurate boundary matching but may fragment structure, whereas semantic and diffusion-based supervision promote visual realism while increasing primitive complexity. Hybrid approaches that combine semantic guidance with constrained representations appear to balance these effects more effectively, suggesting that supervision design plays a central role in regulating structural outcomes.

A recurring insight across model families is the presence of a trade-off between visual realism and editability. Methods optimized for photorealistic appearance often introduce structural complexity that reduces usability, while approaches that enforce layered or constrained representations tend to produce outputs that are easier to modify. This trade-off highlights that SVG generation quality cannot be evaluated using visual metrics alone and must consider structural properties that affect downstream design workflows.

To address this gap, the thesis proposed a multi-dimensional perspective on SVG evaluation. Rather than relying on a single metric, SVG quality should be considered across semantic alignment, structural cleanliness, editability, geometric stability, and computational practicality. This perspective clarifies why methods with similar visual performance may differ substantially in usability and provides a consistent lens for comparing approaches across representation and supervision choices.

Based on the synthesis of findings, several conditions were identified as recurring contributors to clean and editable SVG generation. These include constraints on primitive count, decomposition aligned with semantic structure, supervision applied within structured representation spaces, informed initialization strategies, and architectural priors that reflect design logic. Importantly, these conditions appear across different model categories, suggesting that editability emerges from coordinated design decisions rather than a single methodological paradigm.

Overall, this thesis contributes an analytical understanding of how SVG generation methods construct vector graphics and why structural outcomes differ across approaches. By framing SVG generation as a structured design problem rather than purely an image synthesis task, the work highlights the importance of aligning representation, supervision, and evaluation with the properties that make vector graphics useful in practice.

Future research may explore evaluation protocols that explicitly measure editability and structural quality, develop hybrid architectures that combine fast generation with controllable refinement, and integrate designer-oriented priors into generative pipelines. Advancing SVG generation will likely depend not only on improving visual realism but on developing models that produce vector graphics that remain interpretable, efficient, and practically usable.